\documentclass[12pt]{article}

\usepackage[utf8]{inputenc}
\usepackage[T1]{fontenc}
\usepackage[spanish]{babel}
\usepackage{amsmath}
\usepackage{amssymb}
\usepackage{tikz}
\usepackage{xcolor}
\usepackage{geometry}
\usepackage{enumitem}

% ------------------------------------------------
% MÁRGENES
% ------------------------------------------------
\geometry{
    a4paper,
    left=1.5cm,
    right=1.5cm,
    top=1.3cm,
    bottom=1.5cm
}

% Sin sangría al comenzar párrafos
\setlength{\parindent}{0pt}

% Elimina encabezado y número de página
\pagestyle{empty}

\begin{document}

\noindent
\textbf{1. \textcolor{red}{(20 pts)}} 
En el plano $xy$ se tiene un arco de circunferencia de radio $R$, 
centrado en el origen $O$ y densidad lineal de carga uniforme $\lambda_0$. A partir del punto del arco que interscta con el eje x, se encuentra una barra recta de longitud $2R$, cuya densidad lineal de carga es uniforme e igual a $\lambda_1$.

\vspace{0.25cm}

\noindent
\begin{minipage}[t]{0.58\textwidth}
\vspace{0pt}

\textbf{Determine:}

\vspace{0.25cm}

\begin{enumerate}[
    label=\alph*),
    leftmargin=0.55cm,
    labelsep=0.15cm,
    itemsep=0.55cm,
    topsep=0pt
]

    \item El potencial eléctrico total en el origen $O$ debido a 
    ambas distribuciones de carga. \textcolor{red}{\textbf{(12 pts)}}

    \item El valor de $\lambda_0$, en función de $\lambda_1$, para 
    que el potencial eléctrico total en el origen sea nulo. \textcolor{red}{\textbf{(3 pts)}}

    \item Utilizando el valor de $\lambda_0$ obtenido en el apartado 
    anterior, determine la carga eléctrica total producida por ambas distribuciones. \textcolor{red}{\textbf{(5 pts)}}

\end{enumerate}

\end{minipage}
\hfill
%
\begin{minipage}[t]{0.40\textwidth}
\vspace{-0.35cm}
\centering

\begin{tikzpicture}[scale=1.15]

    % Radio gráfico
    \def\R{1.45}

    % ------------------------------------------------
    % EJES
    % ------------------------------------------------
    \draw[->,thick]
        (-0.45,0) -- (5.05,0)
        node[right] {$x$};

    \draw[->,thick]
        (0,-1.95) -- (0,2.15)
        node[above] {$y$};

    % Origen
    \fill (0,0) circle (1.2pt);
    \node[below left] at (0,0) {$O$};

    % ------------------------------------------------
    % ARCO
    % ------------------------------------------------
    \draw[line width=1.6pt]
        (0,-\R)
        arc[
            start angle=-90,
            end angle=45,
            radius=\R
        ];

    % Radio a 45 grados
    \draw[densely dashed]
        (0,0) -- (45:\R);

    % Ángulo de 45 grados
    \draw
        (0.42,0)
        arc[
            start angle=0,
            end angle=45,
            radius=0.42
        ];

    \node at (0.68,0.23)
        {$45^\circ$};

    % Densidad del arco
    \node[left]
        at (1.0,-0.92)
        {$\lambda_0$};

    % ------------------------------------------------
    % BARRA
    % ------------------------------------------------
    \draw[line width=2.5pt]
        (\R,0) -- ({3*\R},0);

    % Extremos de la barra
    \draw
        (\R,0.11) -- (\R,-0.11);

    \draw
        ({3*\R},0.11) -- ({3*\R},-0.11);

    % Densidad de la barra
    \node[above]
        at ({2.05*\R},0.32)
        {$\lambda_1$};

    % ------------------------------------------------
    % RADIO R
    % ------------------------------------------------
    \node[left]
        at (0,-1.42)
        {$R$};

    % ------------------------------------------------
    % LONGITUD 2R
    % ------------------------------------------------
    \draw[<->]
        (\R,-0.48) -- ({3*\R},-0.48);

    \node[
        fill=white,
        inner sep=1pt
    ]
        at ({2*\R},-0.48)
        {$2R$};

\end{tikzpicture}

\end{minipage}





\newpage





\noindent
\textbf{2. \textcolor{red}{(20 pts)}} 
En el plano $xy$ se tiene un arco de circunferencia de radio $R$, centrado en el origen $O$, el cual posee una densidad lineal de carga uniforme $\lambda_1$. A partir del punto del arco ubicado sobre el eje $x$ negativo se extiende una barra recta de longitud $2R$, con densidad lineal 
de carga uniforme $\lambda_0$. 

\vspace{0.25cm}

\noindent
\begin{minipage}[t]{0.58\textwidth}
\vspace{0pt}

\textbf{Determine:}

\vspace{0.25cm}

\begin{enumerate}[
    label=\alph*),
    leftmargin=0.55cm,
    labelsep=0.15cm,
    itemsep=0.55cm,
    topsep=0pt
]

    \item El potencial eléctrico total en el origen $O$ debido a 
    ambas distribuciones de carga. 
    \textcolor{red}{\textbf{(12 pts)}}

    \item El valor de $\lambda_0$, en función de $\lambda_1$, para 
    que el potencial eléctrico total en el origen sea nulo. 
    \textcolor{red}{\textbf{(3 pts)}}

    \item Utilizando el valor de $\lambda_0$ obtenido en el apartado 
    anterior, determine la carga eléctrica total producida por ambas distribuciones. 
    \textcolor{red}{\textbf{(5 pts)}}

\end{enumerate}

\end{minipage}
\hfill
%
\begin{minipage}[t]{0.40\textwidth}
\vspace{-0.45cm}
\centering

\begin{tikzpicture}[scale=1.12]

    % Radio gráfico
    \def\R{1.45}

    % ------------------------------------------------
    % EJES
    % ------------------------------------------------
    \draw[->,thick]
        (-5.15,0) -- (1.45,0)
        node[right] {$x$};

    \draw[->,thick]
        (0,-2.05) -- (0,2.15)
        node[above] {$y$};

    % Origen
    \fill (0,0) circle (1.2pt);
    \node[below right] at (0,0) {$O$};

    % ------------------------------------------------
    % ARCO
    % theta = 3pi/4 hasta 5pi/4
    % ------------------------------------------------
    \draw[line width=1.6pt]
        (135:\R)
        arc[
            start angle=135,
            end angle=225,
            radius=\R
        ];

    % Radios a los extremos del arco
    \draw[densely dashed]
        (0,0) -- (135:\R);

    \draw[densely dashed]
        (0,0) -- (225:\R);

    % ------------------------------------------------
    % ÁNGULOS DE 45 GRADOS
    % ------------------------------------------------



    \node at (-0.58,0.25)
        {$45^\circ$};

  
    \node at (-0.58,-0.25)
        {$45^\circ$};

    % ------------------------------------------------
    % DENSIDAD DEL ARCO
    % ------------------------------------------------
    \node[left]
        at (-0.65,1.25)
        {$\lambda_1$};

    % Radio R
    \node[right]
        at (-0.78,-0.95)
        {$R$};

    % ------------------------------------------------
    % BARRA
    % Desde x=-3R hasta x=-R
    % ------------------------------------------------
    \draw[line width=2.5pt]
        ({-3*\R},0) -- (-\R,0);

    % Extremos de la barra
    \draw
        (-\R,0.11) -- (-\R,-0.11);

    \draw
        ({-3*\R},0.11) -- ({-3*\R},-0.11);

    % Densidad de la barra
    \node[above]
        at ({-2*\R},0.20)
        {$\lambda_0$};

    % ------------------------------------------------
    % LONGITUD 2R
    % ------------------------------------------------
    \draw[<->]
        ({-3*\R},-0.48) -- (-\R,-0.48);

    \node[
        fill=white,
        inner sep=1pt
    ]
        at ({-2*\R},-0.48)
        {$2R$};

\end{tikzpicture}

\end{minipage}






\newpage


\section*{Pauta V1}

Considerando

\[
V=k\int\frac{dq}{r}
\]

\subsection*{a) \hfill \textcolor{red}{\textbf{(12 pts)}}}

El potencial en el origen está dado por

\[
V_O=V_A+V_B
\qquad
\textcolor{red}{\textbf{1 pt}}
\]

Para cada caso es necesario encontrar $dq$, $r$ y los límites de integración.

\vspace{0.3cm}

\textbf{Arco de circunferencia}

\[
dq_A=\lambda_0R\,d\theta
\qquad
\textcolor{red}{\textbf{1 pt}}
\]

\[
r_{AO}=R
\qquad
\textcolor{red}{\textbf{1 pt}}
\]

\[
\frac{3\pi}{2}
\longrightarrow
2\pi+\frac{\pi}{4}
=
\frac{9\pi}{4}
\qquad
\textcolor{red}{\textbf{1 pt}}
\]

\vspace{0.3cm}

\textbf{Barra}

\[
dq_B
=
\lambda_1\,dx
\qquad
\textcolor{red}{\textbf{1 pt}}
\]

\[
r_{BO}=x
\qquad
\textcolor{red}{\textbf{1 pt}}
\]

\[
R\longrightarrow3R
\qquad
\textcolor{red}{\textbf{1 pt}}
\]

\vspace{0.4cm}

\textbf{Reemplazando}

Para el arco,

\[
V_A
=
\int_{\frac{3\pi}{2}}^{\frac{9\pi}{4}}
k\frac{dq_A}{r_{AO}}
=
\int_{\frac{3\pi}{2}}^{\frac{9\pi}{4}}
k\frac{\lambda_0R}{R}\,d\theta
\qquad
\textcolor{red}{\textbf{1 pt}}
\]

\[
V_A
=
k\lambda_0
\int_{\frac{3\pi}{2}}^{\frac{9\pi}{4}}
d\theta
=
k\lambda_0
\left(
\frac{9\pi}{4}-\frac{3\pi}{2}
\right)
\]

\[
\boxed{
V_A
=
k\lambda_0\frac{3\pi}{4}
}
\qquad
\textcolor{red}{\textbf{1 pt}}
\]

Para la barra,

\[
V_B
=
\int_R^{3R}
k\frac{\lambda_1}{x}\,dx
=
k\lambda_1
\int_R^{3R}\frac{dx}{x}
\qquad
\textcolor{red}{\textbf{1 pt}}
\]

\[
V_B
=
k\lambda_1
\left[
\ln x
\right]_R^{3R}
=
k\lambda_1
\ln\left(\frac{3R}{R}\right)
\]

\[
\boxed{
V_B=k\lambda_1\ln3
}
\qquad
\textcolor{red}{\textbf{1 pt}}
\]

Por lo tanto,

\[
\boxed{
V_O
=
k\left(
\frac{3\pi}{4}\lambda_0+\lambda_1\ln3
\right)
}
\qquad
\textcolor{red}{\textbf{1 pt}}
\]

% ============================================================

\subsection*{b) \hfill \textcolor{red}{\textbf{(3 pts)}}}

Considerando la condición $V_O=0$,

\[
V_O
=
k\left(
\frac{3\pi}{4}\lambda_0+\lambda_1\ln3
\right)
=
0
\qquad
\textcolor{red}{\textbf{1 pt}}
\]

\[
\frac{3\pi}{4}\lambda_0+\lambda_1\ln3=0
\qquad
\textcolor{red}{\textbf{1 pt}}
\]

Por lo tanto,

\[
\boxed{
\lambda_0
=
-\frac{4\ln3}{3\pi}\lambda_1
}
\qquad
\textcolor{red}{\textbf{1 pt}}
\]

% ============================================================

\subsection*{c) \hfill \textcolor{red}{\textbf{(5 pts)}}}

La carga total está dada por

\[
q_{\text{tot}}
=
q_A+q_B
\qquad
\textcolor{red}{\textbf{1 pt}}
\]

Para el arco,

\[
q_A
=
\lambda_0R
\int_{\frac{3\pi}{2}}^{\frac{9\pi}{4}}d\theta
=
\lambda_0R\frac{3\pi}{4}
\qquad
\textcolor{red}{\textbf{1 pt}}
\]

Utilizando

\[
\lambda_0
=
-\frac{4\ln3}{3\pi}\lambda_1,
\]

se obtiene

\[
q_A
=
\left(
-\frac{4\ln3}{3\pi}\lambda_1
\right)
R\frac{3\pi}{4}
\]

\[
\boxed{
q_A=-\lambda_1R\ln3
}
\qquad
\textcolor{red}{\textbf{1 pt}}
\]

Para la barra,

\[
q_B
=
\lambda_1
\int_R^{3R}dx
=
\lambda_1(3R-R)
\]

\[
\boxed{
q_B=2\lambda_1R
}
\qquad
\textcolor{red}{\textbf{1 pt}}
\]

Finalmente,

\[
q_{\text{tot}}
=
-\lambda_1R\ln3+2\lambda_1R
\]

\[
\boxed{
q_{\text{tot}}
=
\lambda_1R(2-\ln3)
}
\qquad
\textcolor{red}{\textbf{1 pt}}
\]

\newpage

% ============================================================
% PAUTA V2
% ============================================================

\section*{Pauta V2}

Considerando

\[
V=k\int\frac{dq}{r}
\]

\subsection*{a) \hfill \textcolor{red}{\textbf{(12 pts)}}}

El potencial en el origen está dado por

\[
V_O=V_A+V_B
\qquad
\textcolor{red}{\textbf{1 pt}}
\]

Para cada caso es necesario encontrar $dq$, $r$ y los límites de integración.

\vspace{0.3cm}

\textbf{Arco de circunferencia}

\[
dq_A
=
\lambda_1R\,d\theta
\qquad
\textcolor{red}{\textbf{1 pt}}
\]

\[
r_{AO}=R
\qquad
\textcolor{red}{\textbf{1 pt}}
\]

\[
\frac{3\pi}{4}
\longrightarrow
\frac{5\pi}{4}
\qquad
\textcolor{red}{\textbf{1 pt}}
\]

\vspace{0.3cm}

\textbf{Barra}

Aunque la barra se encuentra sobre el eje $x$ negativo, $x$ representa
la distancia entre el origen y cada elemento de carga, por lo que se
considera positiva.

\[
dq_B=\lambda_0\,dx
\qquad
\textcolor{red}{\textbf{1 pt}}
\]

\[
r_{BO}=x
\qquad
\textcolor{red}{\textbf{1 pt}}
\]

\[
R\longrightarrow3R
\qquad
\textcolor{red}{\textbf{1 pt}}
\]

\vspace{0.4cm}

\textbf{Reemplazando}

Para el arco,

\[
V_A
=
\int_{\frac{3\pi}{4}}^{\frac{5\pi}{4}}
k\frac{\lambda_1R}{R}\,d\theta
\qquad
\textcolor{red}{\textbf{1 pt}}
\]

\[
V_A
=
k\lambda_1
\int_{\frac{3\pi}{4}}^{\frac{5\pi}{4}}d\theta
=
k\lambda_1
\left(
\frac{5\pi}{4}-\frac{3\pi}{4}
\right)
\]

\[
\boxed{
V_A=\frac{\pi}{2}k\lambda_1
}
\qquad
\textcolor{red}{\textbf{1 pt}}
\]

Para la barra,

\[
V_B
=
\int_R^{3R}
k\frac{\lambda_0}{x}\,dx
=
k\lambda_0
\int_R^{3R}\frac{dx}{x}
\qquad
\textcolor{red}{\textbf{1 pt}}
\]

\[
V_B
=
k\lambda_0
\left[
\ln x
\right]_R^{3R}
=
k\lambda_0
\ln\left(\frac{3R}{R}\right)
\]

\[
\boxed{
V_B=k\lambda_0\ln3
}
\qquad
\textcolor{red}{\textbf{1 pt}}
\]

Por lo tanto,

\[
\boxed{
V_O
=
k\left(
\frac{\pi}{2}\lambda_1+\lambda_0\ln3
\right)
}
\qquad
\textcolor{red}{\textbf{1 pt}}
\]

% ============================================================

\subsection*{b) \hfill \textcolor{red}{\textbf{(3 pts)}}}

Considerando la condición $V_O=0$,

\[
V_O
=
k\left(
\frac{\pi}{2}\lambda_1+\lambda_0\ln3
\right)
=
0
\qquad
\textcolor{red}{\textbf{1 pt}}
\]

\[
\frac{\pi}{2}\lambda_1+\lambda_0\ln3=0
\qquad
\textcolor{red}{\textbf{1 pt}}
\]

Por lo tanto,

\[
\boxed{
\lambda_0
=
-\frac{\pi}{2\ln3}\lambda_1
}
\qquad
\textcolor{red}{\textbf{1 pt}}
\]

% ============================================================

\subsection*{c) \hfill \textcolor{red}{\textbf{(5 pts)}}}

La carga total está dada por

\[
q_{\text{tot}}
=
q_A+q_B
\qquad
\textcolor{red}{\textbf{1 pt}}
\]

Para el arco,

\[
q_A
=
\lambda_1R
\int_{\frac{3\pi}{4}}^{\frac{5\pi}{4}}d\theta
=
\lambda_1R\frac{\pi}{2}
\]

\[
\boxed{
q_A=\frac{\pi}{2}\lambda_1R
}
\qquad
\textcolor{red}{\textbf{1 pt}}
\]

Para la barra,

\[
q_B
=
\lambda_0
\int_R^{3R}dx
=
\lambda_0(3R-R)
\]

\[
\boxed{
q_B=2\lambda_0R
}
\qquad
\textcolor{red}{\textbf{1 pt}}
\]

Utilizando

\[
\lambda_0
=
-\frac{\pi}{2\ln3}\lambda_1,
\]

se obtiene

\[
\boxed{
q_B
=
-\frac{\pi}{\ln3}\lambda_1R
}
\qquad
\textcolor{red}{\textbf{1 pt}}
\]

Finalmente,

\[
q_{\text{tot}}
=
\frac{\pi}{2}\lambda_1R
-
\frac{\pi}{\ln3}\lambda_1R
\]

\[
\boxed{
q_{\text{tot}}
=
\pi\lambda_1R
\left(
\frac{1}{2}-\frac{1}{\ln3}
\right)
}
\qquad
\textcolor{red}{\textbf{1 pt}}
\]
















\end{document}